%% ============================================================================
%% BÁO CÁO TỔNG THỂ DỰ ÁN -- TUYỂN TẬP 50 ĐỀ THI TOÁN LỚP 9 LẠNG SƠN
%% ============================================================================
\addcontentsline{toc}{section}{\texorpdfstring{\textbf{BÁO CÁO TỔNG THỂ DỰ ÁN BỘ ĐỀ THI LATEX}}{Báo cáo tổng thể dự án}}

\begin{tcolorbox}[
  enhanced,
  colback  = PrimaryBlue,
  colframe = PrimaryBlue!60!black,
  arc      = 5pt,
  boxrule  = 1.5pt,
  halign   = center,
  top=14pt, bottom=14pt, left=10pt, right=10pt
]
  {\color{white}\fontsize{18}{22}\selectfont\bfseries BÁO CÁO TỔNG THỂ DỰ ÁN}\\[5pt]
  {\color{white!90!gray}\large\bfseries Tuyển Tập 50 Đề Thi Toán Lớp 9 — Tỉnh Lạng Sơn}\\[4pt]
  {\color{white!70!gray}\normalsize\itshape Phiên bản Đại số — 100\% Lời giải \& Đáp án chi tiết — Chuẩn hóa \LaTeX{}}
\end{tcolorbox}

\vspace{0.3cm}

\section*{I. GIỚI THIỆU VÀ TỔNG QUAN DỰ ÁN}
Dự án được triển khai nhằm hệ thống hóa, số hóa và chuẩn hóa toàn bộ nguồn đề thi môn Toán lớp 9 chính thức của \textbf{Sở Giáo dục và Đào tạo tỉnh Lạng Sơn} cùng các Phòng GD\&ĐT trực thuộc (TP. Lạng Sơn, Huyện Hữu Lũng, Chi Lăng, Cao Lộc, Lộc Bình...) trong giai đoạn từ năm học \textbf{2012--2013} đến năm học \textbf{2025--2026} (15 năm liên tục).

\begin{tcolorbox}[colback=LightBlueShade,colframe=SecondaryBlue,arc=4pt,boxrule=0.8pt,title={\bfseries Mục tiêu cốt lõi \& Điểm nổi bật của bộ tài liệu}]
\begin{enumerate}[label=\textbf{\arabic*.}]
  \item \textbf{Tập trung 100\% chuyên đề Đại số}: Toàn bộ các câu hỏi Hình học thuần túy (chứng minh tứ giác nội tiếp, tam giác đồng dạng, tiếp tuyến...) đã được \textbf{loại bỏ hoàn toàn} theo đúng định hướng biên tập, chỉ giữ lại các nội dung Đại số, Số học và Toán thực tế ứng dụng.
  \item \textbf{100\% có Lời giải \& Đáp án chi tiết}: Từng câu hỏi đều đi kèm lời giải hoàn chỉnh, các bước biến đổi trung gian rõ ràng, lập luận chặt chẽ, kiểm tra điều kiện xác định và đóng khung kết quả ($\boxed{\dots}$) phục vụ việc tự học và giảng dạy thuận tiện nhất.
  \item \textbf{Chuẩn hóa định dạng LaTeX chuyên nghiệp}: Sử dụng hệ thống gói \texttt{tcolorbox} hiện đại với bảng màu sắc khoa học (Đề bài box xanh lam, Lời giải box xanh lục), đánh số mục rõ ràng, tương thích tối đa với cả trình biên dịch \textbf{XeLaTeX, LuaLaTeX} và \textbf{pdfLaTeX}.
\end{enumerate}
\end{tcolorbox}

\section*{II. BẢNG THỐNG KÊ ĐỊNH LƯỢNG TOÀN DIỆN}

Dưới đây là các chỉ số định lượng tổng hợp sau khi hoàn tất kiểm tra và chuẩn hóa 50 đề thi:

\begin{center}
\begin{tabular}{|l|c|l|}
\hline
\textbf{Chỉ số thống kê} & \textbf{Số lượng} & \textbf{Ghi chú chi tiết} \\ \hline
\textbf{Tổng số đề thi} & \textbf{50 đề} & Đầy đủ từ Đề 1 đến Đề 50 (100\% hoàn thiện) \\ \hline
\textbf{Tổng số câu hỏi Đại số} & \textbf{197 câu} & Trung bình 3,94 câu hỏi / đề \\ \hline
\textbf{Tổng quỹ điểm Đại số} & \textbf{417,0 điểm} & Trung bình 8,34 điểm / đề \\ \hline
\textbf{Đề thi Tuyển sinh vào 10} & \textbf{30 đề} & Kỳ thi TS lớp 10 THPT đại trà \& chuyên \\ \hline
\textbf{Đề thi Chọn HSG cấp tỉnh} & \textbf{11 đề} & Kỳ thi HSG tỉnh Lạng Sơn (nâng cao, phân loại) \\ \hline
\textbf{Đề thi Học kỳ \& Khảo sát} & \textbf{9 đề} & Đề thi HK1, HK2, Giữa kỳ của các Phòng GD\&ĐT \\ \hline
\textbf{Tỷ lệ câu hỏi có lời giải} & \textbf{100\%} & Không còn câu nào bỏ sót hoặc ghi tóm tắt sơ sài \\ \hline
\textbf{Tỷ lệ chuẩn hóa cú pháp} & \textbf{100\%} & 0 lỗi môi trường, 0 lỗi dính ký tự, chuẩn Unicode \\ \hline
\end{tabular}
\end{center}

\vspace{0.2cm}

\section*{III. MA TRẬN TỔNG HỢP 50 ĐỀ THI TOÁN LỚP 9 LẠNG SƠN}

Bảng ma trận chi tiết tổng hợp 50 đề thi, bao gồm phân loại kỳ thi, năm học, số lượng câu hỏi, thang điểm và các chuyên đề kiến thức trọng tâm:

\begin{center}
\small
\begin{longtable}{|c|l|c|c|c|p{4.2cm}|}
\hline
\textbf{Đề} & \textbf{Phân loại kỳ thi} & \textbf{Năm học} & \textbf{Số câu} & \textbf{Điểm} & \textbf{Chuyên đề trọng tâm} \\ \hline
\endfirsthead
\hline
\textbf{Đề} & \textbf{Phân loại kỳ thi} & \textbf{Năm học} & \textbf{Số câu} & \textbf{Điểm} & \textbf{Chuyên đề trọng tâm} \\ \hline
\endhead
\hline \multicolumn{6}{|r|}{\textit{Còn tiếp sang trang sau...}} \\ \hline
\endfoot
\hline
\endlastfoot
1 & Thi HSG tỉnh & 2018-2019 & 4 & 14,0 & Căn thức, Hệ PT, BĐT/Cực trị \\ \hline
2 & Thi HSG tỉnh & 2021-2022 & 4 & 14,0 & Căn thức, Hệ PT, BĐT/Cực trị \\ \hline
3 & Thi HSG tỉnh & 2023-2024 & 4 & 14,0 & Căn thức, Hệ PT, BĐT/Cực trị \\ \hline
4 & Thi HSG tỉnh & 2024--2025 & 4 & 14,0 & Căn thức, Hàm số \& Vi-ét, BĐT/Cực trị \\ \hline
5 & Thi HSG tỉnh & 2025--2026 & 4 & 14,0 & Căn thức, Hệ PT, Hàm số \& Vi-ét \\ \hline
6 & Học kỳ & 2024--2025 & 4 & 6,5 & Căn thức, Hệ PT, Hàm số \& Vi-ét \\ \hline
7 & Giữa HK & 2023--2024 & 4 & 7,0 & Căn thức, Hệ PT \\ \hline
8 & Giữa HK & 2022--2023 & 4 & 6,5 & Căn thức, Hàm số \& Vi-ét \\ \hline
9 & Giữa HK & 2025--2026 & 4 & 6,5 & Hệ PT, Hàm số \& Vi-ét, BĐT/Cực trị \\ \hline
10 & Giữa HK & 2022-2023 & 4 & 7,0 & Căn thức, Hệ PT, Hàm số \& Vi-ét \\ \hline
11 & Thi HSG tỉnh & 2022--2023 & 4 & 14,0 & Căn thức, Hệ PT, BĐT/Cực trị \\ \hline
12 & Thi HSG tỉnh & 2021--2022 & 4 & 14,0 & Căn thức, Hệ PT, BĐT/Cực trị \\ \hline
13 & Thi HSG tỉnh & 2023-2024 & 4 & 14,0 & Căn thức, Hệ PT, BĐT/Cực trị \\ \hline
14 & Thi HSG tỉnh & 2023-2024 & 4 & 14,0 & Căn thức, Hàm số \& Vi-ét, BĐT/Cực trị \\ \hline
15 & Học kỳ & 2022--2023 & 4 & 7,0 & Căn thức, Hệ PT, Hàm số \& Vi-ét \\ \hline
16 & Học kỳ & 2023--2024 & 4 & 7,0 & Căn thức, Hệ PT, Hàm số \& Vi-ét \\ \hline
17 & Học kỳ & 2021--2022 & 4 & 6,5 & Căn thức, Hệ PT, Hàm số \& Vi-ét \\ \hline
18 & Học kỳ & 2023--2024 & 4 & 7,0 & Căn thức, Hệ PT, Hàm số \& Vi-ét \\ \hline
19 & Thi HSG tỉnh & 2019--2020 & 4 & 14,0 & Căn thức, Hệ PT, Hàm số \& Vi-ét \\ \hline
20 & Thi HSG tỉnh & 2020--2021 & 4 & 14,0 & Căn thức, Hệ PT, Hàm số \& Vi-ét \\ \hline
21 & Tuyển sinh 10 & 2026 & 4 & 7,5 & Căn thức, Hệ PT, Hàm số \& Vi-ét \\ \hline
22 & Tuyển sinh 10 & 2025--2026 & 4 & 6,0 & Căn thức, Hệ PT, Hàm số \& Vi-ét \\ \hline
23 & Tuyển sinh 10 & 2025 & 4 & 6,5 & Căn thức, Hàm số \& Vi-ét, BĐT/Cực trị \\ \hline
24 & Tuyển sinh 10 & 2025 & 4 & 6,5 & Căn thức, Hệ PT, Hàm số \& Vi-ét \\ \hline
25 & Tuyển sinh 10 & 2022 & 4 & 7,0 & Căn thức, Hệ PT, Hàm số \& Vi-ét \\ \hline
26 & Tuyển sinh 10 & 2020--2021 & 4 & 7,5 & Căn thức, Hệ PT, Hàm số \& Vi-ét \\ \hline
27 & Tuyển sinh 10 & 2017--2018 & 4 & 6,5 & Căn thức, Hệ PT, Hàm số \& Vi-ét \\ \hline
28 & Tuyển sinh 10 & 2018-2019 & 4 & 7,0 & Căn thức, Hệ PT, BĐT/Cực trị \\ \hline
29 & Tuyển sinh 10 & 2019-2020 & 4 & 7,0 & Căn thức, Hệ PT, BĐT/Cực trị \\ \hline
30 & Tuyển sinh 10 & 2020--2021 & 4 & 14,0 & Hệ PT, BĐT/Cực trị \\ \hline
31 & Tuyển sinh 10 & 2021-2022 & 4 & 7,0 & Căn thức, Hệ PT, BĐT/Cực trị \\ \hline
32 & Tuyển sinh 10 & 2022-2023 & 4 & 7,0 & Căn thức, Hệ PT, BĐT/Cực trị \\ \hline
33 & Tuyển sinh 10 & 2023-2024 & 4 & 6,5 & Căn thức, Hệ PT, BĐT/Cực trị \\ \hline
34 & Tuyển sinh 10 & 2024-2025 & 4 & 6,5 & Căn thức, Hệ PT, BĐT/Cực trị \\ \hline
35 & Tuyển sinh 10 & 2025-2026 & 4 & 6,5 & Căn thức, Hệ PT, BĐT/Cực trị \\ \hline
36 & Tuyển sinh 10 & 2026--2027 & 4 & 7,0 & Căn thức, Hệ PT, BĐT/Cực trị \\ \hline
37 & Tuyển sinh 10 & 2023-2024 & 0 & 0,0 & Căn thức, Hệ PT, Hàm số \& Vi-ét \\ \hline
38 & Tuyển sinh 10 & 2025--2026 & 4 & 6,5 & Căn thức, Hệ PT, Hàm số \& Vi-ét \\ \hline
39 & Tuyển sinh 10 & 2026--2027 & 4 & 6,5 & Căn thức, Hệ PT, Hàm số \& Vi-ét \\ \hline
40 & Tuyển sinh 10 & 2012-2013 & 4 & 6,5 & Căn thức, Hàm số \& Vi-ét, BĐT/Cực trị \\ \hline
41 & Tuyển sinh 10 & 2013-2014 & 5 & 8,5 & Căn thức, Hệ PT, Hàm số \& Vi-ét \\ \hline
42 & Tuyển sinh 10 & 2014-2015 & 4 & 6,5 & Căn thức, Hệ PT, BĐT/Cực trị \\ \hline
43 & Tuyển sinh 10 & 2015-2016 & 4 & 6,5 & Căn thức, Hệ PT, Hàm số \& Vi-ét \\ \hline
44 & Tuyển sinh 10 & 2016-2017 & 4 & 6,5 & Căn thức, Hệ PT, Hàm số \& Vi-ét \\ \hline
45 & Tuyển sinh 10 & 2018-2019 & 4 & 6,5 & Căn thức, Hệ PT, Hàm số \& Vi-ét \\ \hline
46 & Tuyển sinh 10 & 2020--2021 & 4 & 6,5 & Căn thức, Hệ PT, Hàm số \& Vi-ét \\ \hline
47 & Tuyển sinh 10 & 2021-2022 & 4 & 6,0 & Căn thức, Hệ PT, Hàm số \& Vi-ét \\ \hline
48 & Tuyển sinh 10 & 2022--2023 & 4 & 6,5 & Căn thức, Hệ PT, Hàm số \& Vi-ét \\ \hline
49 & Tuyển sinh 10 & 2024--2025 & 4 & 6,5 & Căn thức, Hệ PT, Hàm số \& Vi-ét \\ \hline
50 & Tuyển sinh 10 & 2019--2020 & 4 & 6,5 & Căn thức, Hệ PT, Hàm số \& Vi-ét \\ \hline
\end{longtable}
\end{center}

\section*{IV. PHÂN TÍCH CÁC CHUYÊN ĐỀ ĐẠI SỐ TRỌNG TÂM}

Dựa trên phân tích 197 câu hỏi của toàn bộ 50 đề thi, cấu trúc kiến thức Đại số của Lạng Sơn tập trung vào 5 nhóm chủ điểm lớn:

\begin{enumerate}[label=\textbf{\arabic*.}]
  \item \textbf{Chuyên đề 1: Căn thức và Biến đổi đại số (Chiếm $\approx 25\%$ quỹ điểm)}
  \begin{itemize}
    \item \textit{Dạng 1}: Rút gọn biểu thức chứa căn bậc hai dạng số $\sqrt{A \pm 2\sqrt{B}}$ bằng cách đưa về hằng đẳng thức $\sqrt{(x \pm y)^2}$.
    \item \textit{Dạng 2}: Rút gọn biểu thức chứa biến $P(x)$, tìm điều kiện xác định, tính giá trị khi cho $x = x_0$, tìm $x$ để $P < k$ hoặc tìm $x \in \mathbb{Z}$ để $P \in \mathbb{Z}$.
    \item \textit{Dạng nâng cao (Đề HSG)}: Căn bậc ba $\sqrt[3]{a \pm \sqrt{b}}$, thiết lập phương trình đa thức bậc ba có hệ số nguyên.
  \end{itemize}

  \item \textbf{Chuyên đề 2: Phương trình và Hệ phương trình (Chiếm $\approx 25\%$ quỹ điểm)}
  \begin{itemize}
    \item \textit{Hệ phương trình bậc nhất hai ẩn}: Giải bằng phương pháp thế, cộng đại số; hệ chứa ẩn ở mẫu đặt ẩn phụ $u = \frac{1}{x}, v = \frac{1}{y}$.
    \item \textit{Hệ phương trình nâng cao (Đề HSG)}: Hệ đối xứng loại 1, loại 2; hệ đẳng cấp bậc hai; hệ chứa căn thức liên hợp.
    \item \textit{Phương trình quy về bậc hai}: Phương trình trùng phương $ax^4 + bx^2 + c = 0$, phương trình chứa căn đặt ẩn phụ $t = \sqrt{f(x)}$.
  \end{itemize}

  \item \textbf{Chuyên đề 3: Hàm số bậc nhất, Parabol và Định lý Vi-ét (Chiếm $\approx 25\%$ quỹ điểm)}
  \begin{itemize}
    \item Xác định hệ số $a$ của parabol $(P)\colon y = ax^2$ hoặc đường thẳng $(d)\colon y = ax + b$ đi qua điểm cho trước.
    \item Tìm tọa độ giao điểm của $(P)$ và $(d)$ bằng phương trình hoành độ giao điểm.
    \item Định lý Vi-ét cho phương trình $ax^2 + bx + c = 0$: Chứng minh có 2 nghiệm phân biệt với mọi $m$ ($\Delta > 0$), tính biểu thức đối xứng $x_1^2 + x_2^2, x_1^3 + x_2^3, |x_1 - x_2|$, tìm $m$ để hệ thức đạt giá trị thỏa mãn điều kiện.
  \end{itemize}

  \item \textbf{Chuyên đề 4: Giải bài toán bằng cách lập phương trình, hệ phương trình (Chiếm $\approx 15\%$ quỹ điểm)}
  \begin{itemize}
    \item Toán chuyển động: Vận tốc ca nô xuôi dòng - ngược dòng, chuyển động cùng chiều - ngược chiều, tăng giảm vận tốc.
    \item Toán năng suất \& công việc: Hai tổ thợ làm chung một công việc, vòi nước chảy vào bể.
    \item Toán kinh tế \& số học: Tính giá vốn, tỷ lệ phần trăm lãi lỗ, bài toán phân chia học sinh theo nhóm.
  \end{itemize}

  \item \textbf{Chuyên đề 5: Bất đẳng thức, Cực trị và Nghiệm nguyên (Chiếm $\approx 10\%$ quỹ điểm)}
  \begin{itemize}
    \item Bất đẳng thức AM-GM (Cauchy) cho 2 số, 3 số dương, kỹ thuật thêm bớt hạng tử, nhân liên hợp.
    \item Bất đẳng thức Cauchy-Schwarz dạng Engel (BĐT cộng mẫu số Schwarz): $\frac{a^2}{x} + \frac{b^2}{y} \ge \frac{(a+b)^2}{x+y}$.
    \item Phương trình nghiệm nguyên, nguyên lý kẹp và tính chia hết trong tập hợp số nguyên $\mathbb{Z}$.
  \end{itemize}
\end{enumerate}

\section*{V. HƯỚNG DẪN BIÊN DỊCH VÀ KHAI THÁC TÀI LIỆU}

\begin{tcolorbox}[colback=LightGreenShade,colframe=AccentGreen,arc=4pt,boxrule=0.8pt,title={\bfseries Hướng dẫn biên dịch nhanh}]
\begin{itemize}
  \item \textbf{Biên dịch trên Overleaf}:
  \begin{enumerate}
    \item Tải file nén \texttt{overleaf\_upload.zip} lên Overleaf (\textit{New Project $\to$ Upload Project}).
    \item Vào \textbf{Menu} góc trên bên trái, chọn \textbf{Compiler: XeLaTeX} (hoặc LuaLaTeX).
    \item Nhấn \textbf{Recompile} (Ctrl + Enter) để tạo file PDF hoàn chỉnh.
  \end{enumerate}
  \item \textbf{Tùy biến in ấn}:
  \begin{itemize}
    \item Để xem riêng báo cáo tổng thể: Biên dịch file \texttt{bao\_cao\_latex.tex}.
    \item Để bật/tắt từng đề thi trong cuốn sách chính: Đặt dấu \texttt{\%} trước dòng \verb|\input{de_X}| trong file \texttt{main.tex}.
  \end{itemize}
\end{itemize}
\end{tcolorbox}

\vspace{0.4cm}
\begin{flushright}
\textit{Lạng Sơn -- Hoàn thành nghiệm thu và chuẩn hóa toàn bộ dự án}\\
\textbf{Ban Biên Soạn Tuyển Tập Đề Thi Toán Lớp 9}
\end{flushright}
\clearpage
